\documentclass[11pt,a4paper]{article}

%% PACHETE MATEMATICE
\usepackage{amsmath,amssymb,amsfonts}
\usepackage{amsthm}
\usepackage{mathpazo}
\usepackage{eulervm}
\usepackage{enumerate}
\usepackage{color}
\usepackage{graphicx}
\usepackage[all]{xy}
\usepackage{xcolor}
\usepackage{framed}
\usepackage[unicode]{hyperref}
\setcounter{MaxMatrixCols}{20}
\usepackage{multicol}
% \usepackage[romanian]{babel}


%% ASPECT
\linespread{1.05}
\usepackage[T1]{fontenc}
\usepackage[utf8x]{inputenc}
\usepackage[margin=2cm]{geometry}
% \usepackage[color]{showkeys}
\usepackage{anyfontsize}
\usepackage{titlesec}
\titleformat*{\section}{\large\bfseries}

\usepackage{fancyhdr}
\pagestyle{fancy}
\rhead{\small \color{gray} Notițe de Adrian Manea}
\lhead{\small \color{gray} M1-ACS, 2017---2018, M. Olteanu}


\usepackage[autostyle = false, style = german]{csquotes}
\usepackage[romanian]{babel}
\newcommand\qq{\enquote}

%% COMENZILE MELE
\renewcommand*{\proofname}{Dem.:}
\newcommand\bb{\mathbb}
\newcommand\dr{\mathrm}
\newcommand\kal{\mathcal}
\newcommand\fk{\mathfrak}
\newcommand\op{\oplus}
\newcommand\ot{\otimes}
\newcommand\ds{\displaystyle}
\newcommand\id{\indent}
\newcommand\nid{\noindent}
\newcommand\seq{\subseteq}
\newcommand\sneq{\subsetneq}
\newcommand\speq{\supseteq}
\newcommand\spneq{\supsetneq}
\newcommand\rar{\rightarrow}
\newcommand\Rar{\Rightarrow}
\newcommand\xrar{\xrightarrow}
\newcommand\lar{\leftarrow}
\newcommand\Lar{\Leftarrow}
\newcommand\xlar{\xleftarrow}
\newcommand\lrar{\leftrightarrow}
\newcommand\Lrar{\Leftrightarrow}
\newcommand\ideal{\trianglelefteq}
\newcommand\ai{\text{ a.\^{i}. }}
\newcommand\sk{(\textit{Schi\c{t}\u{a}}) }
\newcommand\rhup{\rightharpoonup}
\newcommand\rhdn{\rightharpoondown}
\newcommand\lhup{\leftharpoonup}
\newcommand\lhdn{\leftharpoondown}
\newcommand\vid{\emptyset}
\newcommand\wh{\widehat}
\newcommand\ol{\overline}
\newcommand\NN{\mathbb{N}}
\newcommand\RR{\mathbb{R}}
\newcommand\ZZ{\mathbb{Z}}
\newcommand\QQ{\mathbb{Q}}
\newcommand\CC{\mathbb{C}}

%% STILURILE MELE
\newtheoremstyle{dotless}{}{}{\itshape}{}{\bfseries}{:}{ }{}

\theoremstyle{dotless}
\newtheorem{theorem}{Teorem\u{a}}[section]
\newtheorem{proposition}{Propozi\c{t}ie}[section]
\newtheorem{lemma}{Lem\u{a}}[section]
\newtheorem{corollary}{Corolar}[section]
\newtheorem{exercise}{Exerci\c{t}iu}

\newtheoremstyle{dotlessdr}{}{}{}{}{\bfseries}{:}{ }{}
\theoremstyle{dotlessdr}
\newtheorem{example}{Exemplu}[section]
\newtheorem{definition}{Defini\c{t}ie}[section]
\newtheorem{remark}{Observa\c{t}ie}[section]




\begin{document}

\begin{center}
  {\large\textbf{Seminar 11 \\ Integrale curbilinii}}
\end{center}

\vspace{1cm}

\section{Drumuri parametrizate}
\label{sec:dr-par}

\begin{definition}\label{def:dr-par}
  Fie $J \seq \RR$. Se numește \emph{drum parametrizat} pe $J$ cu valori în $\RR^n$
  orice aplicație continuă $\gamma : J \to \RR^n$.

  Pe componente, $\gamma(t) = (\gamma_1(t), \gamma_2(t), \dots, \gamma_n(t))$,
  relațiile $x_i = \gamma_i(t)$ se numesc \emph{ecuațiile parametrice} ale drumului.
\end{definition}

Dacă $J$ este un interval $[a,b]$, atunci $\gamma(a)$ și $\gamma(b)$ se numesc
\emph{capetele} drumului, iar drumul se numește \emph{închis} dacă $\gamma(a) = \gamma(b)$.

Dacă $\gamma : [a, b] \to \RR^n$ este un drum, \emph{opusul său} se definește prin:
\[
  \gamma^- : [a, b] \to \RR^n, \quad \gamma^-(t) = \gamma(a + b - t).
\]

Date două drumuri parametrizate, ele se pot \emph{concatena}, rezultatul
fiind un drum care parcurge, pe rînd, intervalele de definiție ale celor
două componente.

\begin{remark}\label{rk:dr-concat}
  Concatenarea a două drumuri parametrizate are sens numai atunci cînd capătul
  de final al unui drum coincide cu capătul de început al celuilalt. Mai precis,
  dacă avem: $\gamma_1 : [a, b] \to \RR^n$ și $\gamma_2: [b, c] \to \RR^n$, atunci
  \[
    \gamma_1 \cup \gamma_2 =
    \begin{cases}
      \gamma_1(t), & t \in [a, b] \\
      \gamma_2(t), & t \in [b, c].
    \end{cases}
  \]
\end{remark}

Aspectele analitice de interes sînt:

\begin{definition}\label{def:neted}
  Un drum $\gamma : J \to \RR^n$ se numește \emph{neted} dacă aplicația $\gamma$
  este de clasă $\kal{C}^1(J)$, iar $\gamma'(t) \neq 0, \forall t \in J$.

  Un drum se numește \emph{neted pe porțiuni} dacă se obține prin concatenarea
  unui număr finit de drumuri netede.
\end{definition}

În cazurile particulare de interes, adică în $\RR^2$ și $\RR^3$, vom
nota drumurile parametrizate prin $\gamma(t) = (x(t), y(t))$, respectiv
$\gamma(t) = (x(t), y(t), z(t))$.

%%%%%%%%%%%%%%%%%%%%%%%%%%%%%%%%%%%%%%%%%%%%%%%%%%%%%%%%%%%%%%%%%%%%%% 

\section{Integrala curbilinie}
\label{sec:int-curb}

\textbf{Integrala curbilinie de prima speță}

Dat un drum neted $\gamma : [a, b] \to \RR^3$, iar $f : D \to \RR$, o funcție
continuă cu $D \supseteq \gamma([a, b])$, se definește integrala curbilinie
de prima speță a funcției $f$ pe drumul $\gamma$ prin:
\[
  \int_\gamma f(x, y, z) ds = \int_a^b f(x(t), y(t), z(t)) \cdot %
  \sqrt{(x'(t))^2 + (y'(t))^2 + (z'(t))^2} dt.
\]

În particular, pentru $f = 1$, obținem $L(\gamma)$, care este
\emph{lungimea drumului $\gamma$}, definit prin:
\[
  L(\gamma) = \int_a^b \sqrt{(x'(t))^2 + (y'(t))^2 + (z'(t))^2} dt.
\]

Un alt caz particular este dat de o interpretare fizică. Presupunem că imaginea
lui $\gamma$ reprezintă un fir material, iar $f$ este funcția de densitate a firului.
Atunci putem calcula masa $M$ și coordonatele centrului de greutate
($x_i^G$) ale firului prin:
\begin{align*}
  M &= \int_\gamma f ds \\
  x_i^G &= \frac{1}{M} \int_\gamma x_i f ds.
\end{align*}

\vspace{1cm}

\textbf{Integrala curbilinie de speța a doua}

Fie $\alpha = P dx + Q dy + R dz$ o 1-formă diferențială, unde $P, Q, R$ sînt
funcții continue definite pe $D \seq \RR^3$, mulțime deschisă.

Fie $\gamma : [a, b] \to \RR^3, \gamma(t) = (x(t), y(t), z(t))$ un drum
parametrizat neted, cu $\gamma([a, b]) \seq D$.

Integrala curbilinie a formei diferențiale $\alpha$ de-a lungul drumului $\gamma$
se definește prin:
\[
  \int_\gamma \alpha = \int_a^b (P(\gamma(t))x'(t) + Q(\gamma(t))y'(t) + %
  R(\gamma(t))z'(t) dt.
\]

Putem da și în acest caz o interpretare fizică, folosind \emph{notații vectoriale}.
Asociem 1-formei diferențiale $\alpha = P dx + Q dy + R dz$ un cîmp de vectori
$\vec{V} : D \to \RR^3$, de componente $\vec{V} = (P, Q, R)$. Integrala 1-formei
$\alpha$ pe drumul $\gamma$ devine, în acest caz, \emph{circulația cîmpului} $\vec{V}$
de-a lungul drumului $\gamma$, notată $\int_\gamma \vec{V} \cdot d\vec{r}$.

În particular, dacă interpretăm $\vec{V} = \vec{F}$ ca un cîmp de forțe, atunci
circulația $\int_\gamma \vec{F} \cdot d \vec{r}$ este \emph{lucrul mecanic} efectuat
de forța $\vec{F}$, acționînd pe drumul $\gamma$.

%%%%%%%%%%%%%%%%%%%%%%%%%%%%%%%%%%%%%%%%%%%%%%%%%%%%%%%%%%%%%%%%%%%%%% 

\section{Forme diferențiale exacte}
\label{sec:exact}

\begin{definition}\label{def:exact}
  O 1-formă diferențială $\alpha = P dx + Q dy + R dz$ se numește \emph{exactă} pe o
  mulțime $D$ dacă există o funcție $f$, numită \emph{potențial scalar} sau
  \emph{primitivă}, de clasă $\kal{C}^1(D)$, astfel încît $Df = \alpha$ sau,
  pe componente:
  \[
    \frac{\partial f}{\partial x} = P, \frac{\partial f}{\partial y} = Q, %
    \frac{\partial f}{\partial z} = R.
  \]
\end{definition}

În interpretarea vectorială, cîmpul de vectori $\vec{V} = (P, Q, R)$ asociat formei
diferențiale $\alpha$ se numește \emph{cîmp de gradienți}, deoarece $V = \nabla f$.

\begin{definition}\label{def:inchis}
  O 1-formă diferențială $\alpha = P dx + Q dy + R dz$ se numește \emph{închisă} pe $D$
  dacă sînt verificate, în orice punct din $D$, egalitățile:
  \[
    \frac{\partial P}{\partial y} = \frac{\partial Q}{\partial x}, %
    \frac{\partial Q}{\partial z} = \frac{\partial R}{\partial y}, %
    \frac{\partial R}{\partial x} = \frac{\partial P}{\partial z}.
  \]
\end{definition}

Exactitatea formelor diferențiale ne permite să calculăm mai simplu unele
integrale curbilinii, după cum arată următoarea teoremă:

\begin{theorem}\label{thm:indep-drum}
  Fie $\alpha = Df$ o 1-formă diferențială exactă pe $D$ și fie $\gamma$ un drum
  parametrizat neted, cu imaginea inclusă în $D$, de capete $p, q \in D$. Atunci:
  \begin{enumerate}[(a)]
  \item $\ds\int_\gamma Df = f(q) - f(p)$;
  \item Dacă $\gamma$ este închis, atunci $\int_\gamma Df = 0$.
  \end{enumerate}
\end{theorem}

Să remarcăm faptul că, folosind teorema de simetrie a lui Schwartz, rezultă
că \emph{orice formă diferențială exactă este și închisă}. Reciproca este,
în general falsă. De exemplu, forma diferențială:
\[
  \alpha = \frac{-y}{x^2 + y^2} dx + \frac{x}{x^2 + y^2} dy
\]
este închisă pe $\RR^2 - \{(0, 0)\}$, dar nu este exactă.

Un rezultat fundamental este următorul:

\begin{theorem}[Poincar\'{e}]\label{thm:poinc}
  Fie $\alpha$ o 1-formă diferențială de clasă $\kal{C}^1$, închisă pe domeniul
  $D \seq \RR^n$. Atunci, pentru orice $x \in D$, există o vecinătate deschisă a
  sa $U \seq D$ și o funcție $f$ de clasă $\kal{C}^1$, astfel încît $Df = \alpha$,
  \emph{pe $U$}.
\end{theorem}

Cu alte cuvinte, orice 1-formă diferențială închisă este \emph{local exactă}.

Mai avem nevoie și de următoarea:

\begin{definition}\label{def:stelat}
  O mulțime $D \seq \RR^n$ se numește \emph{stelată} (\emph{domeniu stelat}) dacă
  există un punct $x_0 \in D$, astfel încît segmentul $[x_0, x] \seq D$, pentru orice
  $x \in D$.
\end{definition}

%%%%%%%%%%%%%%%%%%%%%%%%%%%%%%%%%%%%%%%%%%%%%%%%%%%%%%%%%%%%%%%%%%%%%% 

\section{Exerciții}
\label{sec:exc}

1. Să se calculeze următoarele integrale curbilinii de speța întîi:
\begin{enumerate}[(a)]
\item $\int_C x ds$, unde $C: y = x^2, x \in [0, 2]$;
\item $\int_C y^5 ds$, unde $C: x = \frac{y^4}{4}, y \in [0, 2]$;
\item $\int_C x^2 ds$, unde $C: x^2 + y^2 = 2, x, y \geq 0$;
\item $\int_C yds$, unde $C: x(t) = \ln(\sin t) - \sin^2 t$,
  $y(t) = \frac{1}{2} \sin 2t$;
\item $\int_C xyds$, unde $C: x(t) = |t|, y(t) = \sqrt{1 - t^2}$, $t \in [-1, 1]$;
\item $\int_C |x - y| ds$, unde $C: x(t) = |\cos t|, y(t) = \sin t$, $t \in [0, \pi]$.
\end{enumerate}

\vspace{1cm}

2. Să se calculeze următoarele integrale curbilinii de speța a doua:
\begin{enumerate}[(a)]
\item $\int_C \frac{dx}{2a + y} - \frac{dy}{a + x}$, unde $C$ este o curbă
  simplă formată din porțiunea din cercul $x^2 + y^2 + 2ay = 0$ ($a > 0$),
  pentru care $x + y \geq 0$, iar capătul de pornire este $A(a, -a)$;
\item $\int_C xdy - ydx$, unde $C: x^2 + y^2 = 4, y = x \sqrt{3}, x \geq 0, %
  y = \frac{x}{\sqrt{3}}$;
\item $\int_\gamma (x + y) dx + (x - y) dy$, unde domeniul este:
  \[
    \gamma = \{(x, y) \mid x^2 + y^2 = 4, y \geq 0\}.
  \]
\item $\int_\gamma \frac{y}{x + 1} dx + dy$, unde $\gamma$ este triunghiul cu
  vîrfurile $A(2, 0), B(0, 0), C(0, 2)$;
\item $\int_\gamma xdy - ydx$, unde:
  \[
    \gamma = \{(x, y) \mid \frac{x^2}{a^2} + \frac{y^2}{b^2} = 1 \}.
  \]
\end{enumerate}

\vspace{1cm}

3. Să se calculeze $\int_\gamma ydx + xdy$, pe un drum de la $A(2, 1)$
la $B(1, 3)$.

\vspace{1cm}

4. Fie $a \in \RR$ și $P, Q : \RR^2 \to \RR$, definite prin:
\[
  P(x, y) = x^2 + 6y, \quad Q(x, y) = 3ax - 4y.
\]

Să se afle $a$ astfel încît $\omega = Pdx + Qdy$ să fie o 1-formă diferențială
exactă pe $\RR^2$ și să se determine primitive sa ($f \in \kal{C}^1(\RR^2)$, cu
$Df = \omega$).

\vspace{1cm}

5. Să se calculeze circulația cîmpului de vectori $\vec{V}$ de-a lungul curbei
$\gamma$, pentru:
\begin{enumerate}[(a)]
\item $\vec{V} = -(x^2 + y^2) \vec{i} - (x^2 - y^2) \vec{j}$, cu
  \[
    \gamma = \{(x, y) \mid x^2 + y^2 = 4, y < 0 \} \cup %
    \{(x, y) \mid x^2 + y^2 - 2x = 0, y \geq 0 \}.
  \]
\item $\vec{V} = x\vec{i} + xy \vec{j} + xyz \vec{k}$, unde:
  \[
    \gamma = \{(x, y, z) \mid x^2 + y^2 = 1 \} \cap
    \{(x, y, z) \mid x + z = 3 \}.
  \]
\end{enumerate}

\vspace{1cm}

6. Să se calculeze coordonatele centrului de greutate al unui fir cu forma unui
arc de cerc, de rază $R$ și de măsură $\alpha \in (0, \pi)$, presupus omogen.

\vspace{1cm}

7. Să se calculeze masa firului material $\gamma$, cu ecuațiile parametrice:
\[
  \begin{cases}
    x(t) &= t \\
    y(t) &= \frac{1}{2} t^2 \\
    z(t) &= \frac{1}{3} t^3
  \end{cases}, t \in [0, 1]
\]
și cu densitatea $F(x, y, z) = \sqrt{2y}$.

\end{document}
